\section{Operation Declaration Adapters}
\label{sec:declaration-adapters}
An operation declaration can use field names different from those in the
mathematical model. An adapter identifies the corresponding model roles
from the declaration's attribute schema, axis bindings, output order, and
value expression. For a contraction, these must satisfy the contract in
Table~\ref{tab:formal-operations}, including the batch and contraction
positions in both operands. Additional constraints must be checked: a
restricted declaration is not an unrestricted instance of the model.

For an admitted declaration, the fields corresponding to model attributes
are distinct and cover all roles. The RHS adapter therefore gives a
bijection between declared attribute assignments and model assignments.
The LHS adapter reads captured values and preserves their sharing.

The preservation argument follows the expression tree. Inputs retain their
shapes and values. At each operation, matching the model contract preserves
attribute legality, ordered output shape, and element computation, assuming
these are preserved for its children. Thus the adapted tree represents the
same computation, and the bijection neither loses nor introduces attribute
solutions. Recognition uses supported structural templates; it does not
require deciding equivalence of arbitrary operation definitions. This
argument concerns declaration adaptation, not a soundness proof of the
full symbolic compiler.
